% Sections 3.1 and 3.2, from lectures/L03/appendix/a_theory.md: one section per numbered section of
% the appendix.

\section{Neural Networks}
\label{c3:sec:nn}\label{c3:app:a}

\subsection{Introduction}
\begin{itemize}
  \item Through machine learning, a given machine can become capable of learning rules without
    having been explicitly programmed to do so.
  \item \textbf{Deep learning} is a subfield of machine learning and involves techniques that
    vaguely resemble how our brains work. The name stems from the fact that the technique was
    originally inspired by how our brains function. In practice, though, brains and neural networks
    work very differently.
  \item Deep learning uses so-called \textbf{neural networks} containing several steps, or layers,
    between input and output, often hundreds, which allows relevant information to be extracted
    from the input while irrelevant information is discarded. This makes deep learning well suited
    to images and similar data containing a large amount of information, most of which is usually
    insignificant.
  \item The more layers used in a neural network, the deeper the model. Information can be thought
    of as passing through and being processed by different filters that distill the information.
    In each layer, a certain type of irrelevant data is filtered out while relevant data is
    retained. Step by step, the input changes from its original form and increasingly contains
    information about the output.
\end{itemize}

\subsection{Theory Behind Training a Neural Network}
\begin{itemize}
  \item For a neural network to predict with low deviation, the network must be trained so that the
    parameters (weights and biases) of its nodes are adjusted to suitable values.
  \item Existing training sets are therefore used to train the network for a certain number of
    epochs. At each epoch, the order of the training sets is randomized so the network doesn't get
    used to any patterns arising from the particular order the training sets happen to be stored
    in.
\end{itemize}
The following steps are performed for each training set, for a simple neural network consisting of
an input layer, a hidden layer, and an output layer.

\subsubsection{1. Feedforward}
\begin{itemize}
  \item The nodes in the input layer form the inputs $x$ to every node in the hidden layer.
  \item Since these inputs $x$ carry different degrees of importance, also called weight, each is
    multiplied by its own weight $w$ in the hidden layer. These weights should be randomly chosen at
    the start, and then adjusted during training.
\end{itemize}
For each node in the hidden layer, the contribution $w \cdot x$ from each input $x$ from the input
layer is summed together with the node's resting point, or bias, $b$ into a sum $s$:
\[
  s = b + \sum_{i=0}^{n-1} w_i \cdot x_i
    = b + w_0 \cdot x_0 + w_1 \cdot x_1 + \dots + w_{n-1} \cdot x_{n-1},
\]
where:
\begin{itemize}
  \item $b$ is the node's bias.
  \item $n$ is the number of inputs.
  \item $w_i \cdot x_i$ is the product of each weight $w$ and its connected input $x$.
\end{itemize}
The node's predicted output $\yp$ is obtained via an activation function $\sigma$:
\[
  \yp = \sigma(s) = \sigma\!\left(b + \sum_{i=0}^{n-1} w_i \cdot x_i\right)
\]
The outputs from the hidden layer then form the inputs to the output layer, where the same
principle applies.

\subsubsection{2. Backpropagation}
The deviation for each node's output is computed starting from the output layer and moving backward
to the hidden layer. For a given node $n$ in the output layer:
\[
  \delta = \yref - \yp
\]
The computed error $\Delta e$ for a given node $n$ in the output layer:
\[
  \Delta e = \delta \cdot \yp',
\]
where $\yp'$ is the derivative of the node's predicted output.

For a given node $n$ in a given hidden layer, the deviation $\delta$ is computed as
\[
  \delta = \sum_{i=0}^{n-1} \left[\Delta e_i \cdot w_i\right]
\]
and the computed error $\Delta e$ for the node:
\[
  \Delta e = \delta \cdot \yp'
\]

\subsubsection{3. Optimization}
Once the error for the nodes' parameters has been computed, optimization is performed via a gradient
algorithm. The rate of change $\Delta c_n$ for a given node $n$:
\[
  \Delta c_n = \Delta e_n \cdot L,
\]
where:
\begin{itemize}
  \item $\Delta e_n$ is the computed error for the current node.
  \item $L$ is the learning rate.
\end{itemize}
The node's bias $b_n$ is updated:
\[
  b_n = b_n + \Delta c_n
\]
Each weight $w_j$ connected to node $n$ is updated:
\[
  w_j = w_j + \Delta c_n \cdot y_j,
\]
where $y_j$ is the output from the connected node $j$ in the previous layer (for the input layer,
$y_j = x_j$).

\subsection{Activation Functions}
\begin{itemize}
  \item Some type of activation function is typically used for each node in the hidden and output
    layers. This introduces non-linearity into the network, which is required to model anything
    beyond simple linear relationships.
  \item Some activation functions (e.g.\ sigmoid, tanh, and softmax) also keep the output within a
    certain range, while others (e.g.\ ReLU) are unbounded above.
\end{itemize}

\begin{booktable}{@{}>{\raggedright\arraybackslash}p{5.6em}>{\raggedright\arraybackslash}p{6.6em}LL@{}}
  \thead{Activation function} & \thead{Output range} & \thead{Typical use} &
    \thead{Main drawback} \\
  \midrule
  ReLU & $[0, \infty)$ & Hidden layers (default choice) & Can ``die'' on negative inputs \\
  Leaky ReLU & $(-\infty, \infty)$ & Hidden layers, alternative to ReLU & Extra hyperparameter
    ($\alpha$) to choose \\
  Sigmoid & $(0, 1)$ & Output layer for binary classification & Saturates at large $\lvert s
    \rvert$, gives small gradients \\
  Tanh & $(-1, 1)$ & Hidden layers, zero-centered output & Saturates at large $\lvert s \rvert$,
    gives small gradients \\
  Softmax & $(0, 1)$, sums to 1 & Output layer for multi-class classification & Requires one output
    node per class \\
\end{booktable}

\noindent The node's output $y$ is obtained as:
\[
  y = \sigma(s) = \sigma\!\left(b + \sum_{i=0}^{n-1} w_i \cdot x_i\right)
\]

\subsubsection{1. ReLU}
The most common activation function in deep learning is \textbf{ReLU} (\emph{Rectified Linear
Unit}):
\[
  \begin{cases} s > 0 \Rightarrow y = s \\ s \leq 0 \Rightarrow y = 0 \end{cases}
\]
As a rule of thumb, ReLU tends to be the most suitable activation function to implement when
there's no clear-cut answer.

\textbf{Leaky ReLU} is an improved variant where signals below zero are retained as a weak signal:
\[
  \begin{cases} s > 0 \Rightarrow y = s \\ s \leq 0 \Rightarrow y = s \cdot \alpha \end{cases}
\]
\begin{itemize}
  \item As a rule of thumb, $\alpha = 0.01$ can be used.
  \item Leaky ReLU counteracts nodes ``dying'' (getting stuck at zero), which can otherwise cause
    the model to train poorly.
\end{itemize}
The derivative of ReLU:
\[
  \begin{cases} s > 0 \Rightarrow dy = 1 \\ s \leq 0 \Rightarrow dy = 0 \end{cases}
\]
Every derivative below is a function of the weighted sum $s$, which is what the implementation
passes in. Sigmoid and tanh simply happen to have compact forms that reuse the output $y$; that's
the same number, not a different argument. See \secref{c5:sec:math} for why the layer keeps $s$
around alongside $y$.

\subsubsection{2. Sigmoid}
Sigmoid is a common activation function when only two classes are used, where the output is set to
a number between 0 and 1:
\[
  y = \frac{1}{1 + e^{-s}}
\]
The derivative:
\[
  \sigma'(s) = y \cdot (1 - y), \quad \text{where } y = \sigma(s)
\]

\subsubsection{3. Tanh}
Tanh (\emph{hyperbolic tangent}) resembles sigmoid but produces an output between $-1$ and 1, which
often leads to faster learning since the outputs become zero-centered:
\[
  y = \tanh(s) = \frac{e^s - e^{-s}}{e^s + e^{-s}}
\]
The derivative:
\[
  \sigma'(s) = 1 - \tanh^2(s) = 1 - y^2
\]

\subsubsection{4. Softmax}
Softmax is a common activation function in the output layer, where the sum of all output
probabilities equals one. The output $y$ for a node in a layer of $n$ nodes:
\[
  y = \frac{e^s}{\sum_{i=0}^{n-1} e^{s_i}}
\]
Softmax is particularly useful for multi-class classification, e.g.\ image recognition.

% ------------------------------------------------------------------------------------------
\section{Dense Layer Architecture}
\label{c3:sec:dense}

\subsection{Overview}
Starting in this chapter, you'll build up a small neural network in C++ step by step, over this
chapter and the next. The first piece is a \code{dense_layer::Interface} and a placeholder
implementation of it, \code{dense_layer::Stub}:

\begin{console}
ml::dense_layer::Interface   (Interface for a dense layer)
└── ml::dense_layer::Stub    (Dense layer stub)
\end{console}

\noindent In Chapter~\ref{c4:ch:network}, this interface will be used by a
\code{neural_network::Shallow} class to build a small working (if not yet properly trained) network.
Chapter~\ref{c5:ch:dense} then replaces \code{Stub} with a real, trainable implementation; the
interface stays the same throughout, so nothing that depends on it needs to change when that
happens.

\subsection{The Dense Layer's Interface}
\code{ml::dense_layer::Interface} defines the contract every dense layer must fulfill:

\begin{booktable}{@{}>{\raggedright\arraybackslash}p{14em}L@{}}
  \thead{Method} & \thead{Description} \\
  \midrule
  \code{nodeCount()} & Number of nodes in the layer \\
  \code{weightCount()} & Number of weights per node \\
  \code{output()} & The layer's outputs after feedforward \\
  \code{error()} & The layer's computed error after backpropagation \\
  \code{weights()} & The layer's weights (2D: node $\times$ weight) \\
  \code{feedforward(input)} & Computes outputs from the input \\
  \code{backpropagate(reference)} & Computes error from reference values (output layer) \\
  \code{backpropagate(nextLayer)} & Computes error from the next layer (hidden layer) \\
  \code{optimize(input, learningRate)} & Updates bias and weights \\
  \code{initParams()} & Resets the layer's trainable parameters, i.e.\ bias and weights \\
\end{booktable}

Every method in the table is pure virtual: the interface states what a dense layer must be able to
do and nothing about how. \code{initParams()} is what lets the network in
Chapter~\ref{c4:ch:network} start each training run from freshly drawn parameters without knowing
which kind of layer it holds; a layer with nothing to reset, like the stub, implements it as an
empty method. See \secref{c3:sec:exercises}.

\subsection{The Stub Class}
\code{ml::dense_layer::Stub} implements the interface but doesn't do anything meaningful:
\begin{itemize}
  \item Feedforward always sets the output to a fixed value (e.g.\ \code{0.5}).
  \item Backpropagation and optimization do nothing.
  \item The parameter reset is implemented as an empty method: there's no bias, and the weights stay
    at zero.
\end{itemize}
The stub exists so that other code, namely the \code{Shallow} network class you'll build in
Chapter~\ref{c4:ch:network}, can be written, compiled, and test-run against a real
\code{dense_layer::Interface} before a correct implementation exists. A correct implementation
follows in Chapter~\ref{c5:ch:dense}.
